% Source fingerprint: 8b3ba2252dface305b9d0eb6a7dbb94ed4689bc53f0732153c7a0d8ad3228a6c
% T02: raw TeX cells preserved; no numeric highlighting.
\begin{table}[htbp]
\centering\small\renewcommand{\arraystretch}{1.18}\setlength{\tabcolsep}{4pt}
\caption[延拓的存在、唯一性与完备性]{延拓的存在、唯一性与完备性。存在性、唯一性和完备性属于不同步骤。限制可测域可能失去完备性，不能把完整 Carathéodory 可测域与生成代数混为一层。}
\label{b1:tab:measure-extension-objects}
\begin{tabularx}{\linewidth}{@{}>{\raggedright\arraybackslash}p{3.1cm}>{\raggedright\arraybackslash}p{4.7cm}>{\raggedright\arraybackslash}X@{}}
\toprule
对象与定义域 & 本步得到的性质 & 额外条件与边界 \\
\midrule
\(p\) 于 \(\mathcal R\) & 原始前测度；已对允许的可数分解可加 & 它的定义域尚不是所有集合 \\
\(p^*\) 于 \(\mathcal P(X)\) & 由可数覆盖代价生成外测度 & 外测度在任意集合上不必可加 \\
\(p^*|_{\mathcal M_{p^*}}\) & Carathéodory 可测域上的完备测度 & 可加性通过切分判据建立，不要求先有唯一性 \\
\(\mu_{\mathrm C}\) 于 \(\sigma(\mathcal R)\) & 原前测度在目标域上的延拓 & 本章用 \(\sigma\)-有限性保证目标域上的唯一性；限制后不自动完备 \\
\bottomrule
\end{tabularx}
\end{table}
